\subsection{Generalized deletion propagation}\label{sec:gdp}

Generalized deletion propagation (GDP) permits simultaneous deletion and
preservation requirements on multiple query outputs, with shared decisions
on source tuples and an objective expressed through view-deletion counts
\citep[Definition 4.1]{makhija2025gdp}. Its ILP framework supports self-joins
and both set and bag semantics; its tractability results concern specified
deletion-propagation variants and query conditions
\citep[Sections 5--6]{makhija2025gdp}.

Appendix~\ref{app:gdp} derives a polynomial-size GDP encoding of each
unit-cost temporal profile entry. Protected compatibility relations compile
strict timing and root incidence; retaining every event identifier in the
query output preserves tuple multiplicity under set semantics. Deletion and
preservation bounds on the same query enforce an exact count. Mandatory
statuses are handled by normalization and singleton preservation queries.
Polynomially many target pairs recover the full profile table. This is a
derived encoding, not a temporal theorem stated in the cited work.

The appendix also gives a weighted witness-indicator ILP with the original
binary cost coefficients. That objective extends the literal cardinality
objective of GDP Definition~4.1; expanding an event into one copy per cost
unit would not justify polynomial input size. Thus the formulation and
full-table output are not the claimed novelty. The contribution claimed in
Theorem~\ref{thm:main} is the structural algorithm with its joint time and
space guarantee. Obtaining that guarantee through GDP requires a solving-time
and parameter-transfer argument beyond the encoding.
